\documentclass[10pt,a4paper]{amsart}
\usepackage[margin=19mm]{geometry}
\usepackage{amsmath,amssymb,mathtools}
\usepackage[scr=rsfs,cal=euler]{mathalfa}
\usepackage{xfrac}
\usepackage[unicode,hidelinks,bookmarks=false]{hyperref}
\newcommand{\ie}{i.e.\ }
\newcommand{\cf}{cf.\ }
\newcommand{\resp}{respectively\ }
\newcommand{\Subsection}{\subsection}
\providecommand{\nobreakdash}{\nobreak\hbox{-}}
\setlength{\emergencystretch}{2em}
\multlinegap0pt
\usepackage{bm}
\usepackage{bbm}
\usepackage{dsfont}
\usepackage{xspace}
\usepackage{cancel}
\usepackage{mathtools}
\usepackage{amsthm}
\makeatletter
\@ifundefined{theo}{\newtheorem{theo}{Theo}}{}
\@ifundefined{lemm}{\newtheorem{lemm}[theo]{Lemma}}{}
\makeatother
\def\cF{{\mathcal F}}
\let\pa=\partial
\newcommand \tri {\coloneqq}
\newcommand\R{{\mathbb R}}
\newcommand{\beqs}{\begin{equation*}}
\newcommand{\beq}{\begin{equation}}
\newcommand{\eeqs}{\end{equation*}}
\newcommand{\eeq}{\end{equation}}
\title[Transition threshold for the Navier-Stokes-Coriolis system a]{Transition threshold for the Navier-Stokes-Coriolis system at high Reynolds numbers}
\author{Minling Li, Changzhen Sun, Chao Wang, Dongyi Wei, Zhifei Zhang}
\date{Source: arXiv:2511.21294v2 (CC BY 4.0)}
\begin{document}
\maketitle

\noindent\emph{Source and licence.} Transition threshold for the Navier-Stokes-Coriolis system at high Reynolds numbers. Version arXiv:2511.21294v2 (CC BY 4.0), distributed under the Creative Commons Attribution 4.0 International licence, as recorded in the arXiv licence metadata for this exact version and shown on the abstract page https://arxiv.org/abs/2511.21294v2. Full rights notice: licenses/arxiv-2511.21294-CC-BY-4.0.txt.\par\smallskip

\noindent\textit{Editorial scope.} Counted unit: the Lemm whose source label is \texttt{App-dis-1} in the article (source0.tex lines 3442--3454, source label \texttt{App-dis-1}), delivered together with the article's own complete original proof of it (source0.tex lines 3456--3583, one \texttt{proof} environment of 128 lines). No mathematics is written, repaired or re-derived: statement and proof are the article's own text line for line. Required context retained uncounted: none. Every definition, display and label of those lines is retained unchanged. Omitted from this package and not redistributed: 1--3441, 3455--3455, 3584--4095. Nothing inside a retained excerpt is omitted or altered: every retained body is delivered exactly as the article's lines are written, so no omission receipt of any kind accompanies this package. Citations: the retained excerpts carry no citation command at all, so no bibliography entry of the article is transcribed and no reference is redistributed. Reference adaptation: none was needed and none was made; every reference inside the delivered unit resolves inside it, so no unresolved reference or citation is produced. Printed slips kept literal: no printed word, symbol, hypothesis or number of the counted statement or of its proof was altered. Layout and class: the article's own class is \texttt{article} with packages including graphicx, inputenc, amsmath, amssymb, amsthm, mathrsfs, dsfont, bm, geometry, titlesec, hyperref, color, fancyhdr, tikz; the wrapper is the lane's pinned \texttt{amsart} editorial wrapper at 10pt on A4, which restates the theorem-like environments the retained text needs and adds the article's own macro definitions that the retained text needs (\texttt{\textbackslash R}, \texttt{\textbackslash beq}, \texttt{\textbackslash beqs}, \texttt{\textbackslash cF}, \texttt{\textbackslash eeq}, \texttt{\textbackslash eeqs}, \texttt{\textbackslash pa}, \texttt{\textbackslash tri}), with extra wrapper packages (bm, bbm, dsfont, xspace, cancel, mathtools). The delivered numbering is the unit's own. Assets: no retained span contains a figure environment, a graphics-inclusion command, a PDF-page inclusion, a table or a TikZ picture; no image file is copied into this package, the package declares no asset dependency and no image file is redistributed with it. Licence: version arXiv:2511.21294v2 of this article is distributed under the Creative Commons Attribution 4.0 International licence, as recorded in the arXiv licence metadata for this exact version and shown on the abstract page https://arxiv.org/abs/2511.21294v2. A full rights notice is delivered at licenses/arxiv-2511.21294-CC-BY-4.0.txt. Characters: every retained excerpt is pure ASCII, so the delivered engine, XeLaTeX with its default Latin Modern text font, reports no missing character.
\begin{lemm}\label{App-dis-1}
	Let $\mathcal{R}_3=\pa_z|\nabla_{y,z}|^{-1}$.
    %, and $f(y,z)\in\mathcal{S}(\R^2)$.
	There exists $C>0$ such that for any $\gamma>0, f\in \mathcal{S}(\R^2),$
	\beq\label{dispersive}
	\left\| e^{\gamma t\mathcal{R}_3(D)}f \right\|_{L^\infty(\mathbb{R}^2)} \leq C (1+|\gamma t|)^{-\frac{1}{2}} \|f\|_{W^{2^{+},1}(\mathbb{R}^2)}\,.
	\eeq
	%Moreover, for any $q\in [2,+\infty),$ 
	%\beq\label{dispersive-Lq}
%	\left\| e^{\gamma t\mathcal{R}_3}f \right\|_{L^q(\mathbb{R}^2)} \lesssim_{\iota}  (1+|\gamma t|)^{-\frac{1}{2}(1-\f2q)} \|f\|_{W^{(2^{+})(1-\f2q),q'}(\mathbb{R}^2)}\,.
%	\eeq
	%for any $\iota>0.$
\end{lemm}

\begin{proof}
	%The estimate \eqref{dispersive-Lq} follows from the interpolation between \eqref{dispersive} and $\|e^{\gamma t \mathcal{R}_3}\|_{L^2\rightarrow L^2}\leq 1.$ Let us  detail the proof of \eqref{dispersive}. 
	It suffices to prove the case of $\gamma=1.$
	In fact, to show \eqref{dispersive},  we will prove that
	\beq\label{dispersive-1}
	\left\| e^{t\mathcal{R}_3}f \right\|_{L^\infty(\mathbb{R}^2)} \leq |t|^{-\frac{1}{2}} \|f\|_{W^{2^{+},1}(\mathbb{R}^2)}, \quad \forall~ |t|\geq 1,
	\eeq
	which, combined with the Sobolev embeddings, gives \eqref{dispersive}.
	
	Let $\chi_j=\chi_j(\xi, \eta)=\chi_0(2^{-j}(\xi, \eta))\, (j\in \mathbb{Z})$ be radial, dyadic cut-off functions that are supported on $\{2^{j-1}\leq |(\xi,\eta)|\leq 2^{j+1}\}$ and satisfies 
	\begin{align*}
		\sum_{j\in \mathbb{Z}}\chi_j(\xi,\eta)=1, \quad \text{for}~(\xi,\eta)\neq 0\,.
	\end{align*}
	
	Let us first claim that it suffices to show that 
	\begin{align}\label{disp-show}
		\left\| e^{t\mathcal{R}_3}\chi_0(D)f \right\|_{L^\infty(\mathbb{R}^2)} \lesssim |t|^{-\frac{1}{2}} \|\chi_0(D)f\|_{L^1(\mathbb{R}^2)}.
	\end{align}
	Indeed, it holds that 
	\begin{align*}
		\cF\big[e^{t\mathcal{R}_3}\chi_j(D)f\big](\xi,\eta)&=e^{t \mathcal{R}_3(\xi,\eta)}\chi_0\big(2^{-j} (\xi,\eta)\big) \cF[f](\xi,\eta) \\
		&= e^{t \mathcal{R}_3(\tilde{\xi},\tilde{\eta})}\chi_0(\tilde{\xi},\tilde{\eta}) \cF[f](2^{j}(\tilde{\xi},\tilde{\eta}))
        \\
        &=e^{t \mathcal{R}_3(\tilde{\xi},\tilde{\eta})}\chi_0(\tilde{\xi},\tilde{\eta}) \cF[g](\tilde{\xi},\tilde{\eta}),
        %\\
		%&=\cF\big(e^{it P(D)} (\chi_j(D)f)(2^{-j}\cdot) \big)(\tilde{\xi}, \tilde{\eta}),
	\end{align*}
	where we denoted $(\tilde{\xi}, \tilde{\eta})=2^{-j}(\xi,\eta)$ and $\cF[g](\tilde{\xi},\tilde{\eta})\tri \cF[f](2^{j}(\tilde{\xi},\tilde{\eta}))$.
	Consequently, it follows from \eqref{disp-show} that
	\begin{align}\label{dispes-pj}
		\left\| e^{t\mathcal{R}_3}\chi_j(D)f \right\|_{L^\infty(\mathbb{R}^2)} 
        &=2^{2j}\left\| e^{t\mathcal{R}_3}\chi_0(D)g \right\|_{L^\infty(\mathbb{R}^2)} 
        \lesssim 2^{2j}|t|^{-\frac{1}{2}} \left\| \chi_0(D)g \right\|_{L^1(\mathbb{R}^2)} 
        \nonumber
        \\
        &=|t|^{-\frac{1}{2}} \left\| \chi_j(D)f(2^{-j}(x,y))\right\|_{L^1(\mathbb{R}^2)} 
        =
        2^{2j}|t|^{-\frac{1}{2}} \|\chi_j(D)f\|_{L^1(\mathbb{R}^2)}, \quad \forall~ j\in \mathbb{Z}\, ,
	\end{align}
    where we have used the fact that
    \begin{align*}
        (\chi_0(D)g) (x,y)=2^{-2j}(\chi_j(D) f)(2^{-j}(x,y)).
    \end{align*}
	Summarizing $j$ over $\mathbb{Z},$ we obtain 
	\eqref{dispersive-1}.
	
	We now focus on proving \eqref{disp-show}. By Inverse Fourier transform, it suffices to show that
	\[
	K(t, y,z) = \int_{\mathbb{R}^2} e^{i  (y\,\xi + z\,\eta) + t \mathcal{R}_3(\xi, \eta)} \tilde{\chi}_0(\xi, \eta) \, d\xi \, d\eta
	\]
	satisfies (here $ \tilde{\chi}_0={\chi}_0+{\chi}_1+{\chi}_{-1}$)
	\[
	\| K(t,y,z) \|_{L^\infty_{y,z}} \lesssim t^{-\frac{1}{2}}, \quad \forall~ t\geq 1.
	\]
	It is convenient to write $K$ in the polar coordinate
	\[
	\xi = r \cos \theta, \quad \eta = r \sin \theta \Rightarrow \mathcal{R}_3(\xi, \eta) = i\,\sin \theta\,\,.
	\]
	Let us also set \( y = \rho \cos \varphi, \, z = \rho \sin \varphi \), then $K$ can be written
	\begin{align*}
		K(t,y,z)&=K(t,\rho,\varphi) 
		= \int_0^\infty \int_{-\pi}^{\pi} e^{i t \sin \theta} \, e^{i r  %\omega_{\rho}
			\rho \cos (\theta-\varphi)} A(r) \, r \, dr \, d\theta\\
		&=\int_0^\infty \int_{-\pi}^{0} e^{i t \sin \theta} \, e^{i r  \rho \cos (\theta-\varphi)} A(r) \, r \, dr \, d\theta+\int_0^\infty \int_{0}^{\pi} e^{i t \sin \theta} \, e^{i r  \rho \cos (\theta-\varphi)} A(r) \, r \, dr \, d\theta\\
		&\tri  K_{1}(t, \rho, \varphi)+K_{2}(t,\rho, \varphi),
	\end{align*}
	where 
	\beqs
	%\omega_{\rho}(\theta) = \rho \cos(\theta - \varphi), \qquad 
	\tilde{\chi}_0(\xi,\eta)= A(r) \in  C_c^{\infty}(\mathbb{R}).
	\eeqs
	Define
	\[
	B(z) \tri \int_0^\infty e^{i r  z} A(r) r \, dr\,=\mathcal{F}^{-1}(A(r) r)(z).
	\]
	We will prove that, %the desired estimates for $K_2,$ ie. 
    for any $(\rho,\varphi)\in (0,+\infty)\times [0,2\pi]$
	\beq\label{es-K2}
	|K_2(t, \rho,\varphi)|=\bigg|\int_{0}^{\pi} e^{it \sin\theta}B(\rho \cos (\theta-\varphi))\, d\theta\bigg| \lesssim |t|^{-1/2},
	\eeq
	the proof of $K_1$ is similar. Note that $B=\mathcal{F}^{-1}(A(r) r)$ is a Schwartz function, in particular $B\in W^{1,1}(\R)$,
\begin{align*}
		\int_{-\pi}^{\pi} |\partial_{\theta}B(\rho \cos (\theta-\varphi))|\, d\theta=\int_{-\pi}^{\pi} |\partial_{\theta}B(\rho \cos \theta)|=2\int_{0}^{\pi} |\partial_{\theta}B(\rho \cos \theta)|\, d\theta=2\int_{-\rho}^{\rho}|B'(z)|dz\leq C, 
	\end{align*}
 \if0   To do so, we need to consider different situations depending on the sizes of $\rho$ and $t.$
	
	\textbf{Case 1.} When \( 0\leq \rho \leq 10 \). 
	
	
	In this case, 
	\( B_{\rho}(\theta), \partial_\theta B_{\rho}(\theta) \)  are  bounded functions in \( \theta \), uniformly in \( \rho \in [0,10], \varphi\in [0,2\pi]\): 
	\beq \label{unies-B}
	\sup_{\rho\in[0,10]}\sup_{\varphi\in  [0,2\pi]}\sup_{\theta\in[0,\pi]} \Big(|B_{\rho}(\theta-\varphi)|+|\pa_{\theta} B_{\rho}(\theta-\varphi)|\Big)
	<+\infty\,.
	\eeq\fi
	We can apply the 1-d stationary phase arguments to show \eqref{es-K2}.
	Indeed, $(\sin\theta)'=\cos\theta$ vanishes at 
	$\theta_0=\pi/2,$ but $(\sin\theta)''=-\sin\theta$ does not vanish near $\theta=\theta_0.$ 
	We thus split $K_2=K_{21}+K_{22},$ where 
	\begin{align*}
		K_{21}(t,\rho,\varphi)=\int %{\frac{\pi}{2}-\delta}^{\frac{\pi}{2}+\delta}
		_{I_{\delta}}\, e^{it \sin\theta} B(\rho \cos (\theta-\varphi)) \, d\theta,
        ~~\text{and}
		~~  
        K_{22}(t,\rho,\varphi)=%\bigg(\int_0^{\frac{\pi}{2}-\delta}+\int_{\frac{\pi}{2}+\delta}^{\pi}\bigg)
		\int_{[0,\pi]\backslash I_{\delta}}\, e^{it \sin\theta} B(\rho \cos (\theta-\varphi)) \,d\theta, 
	\end{align*}
	where $I_{\delta}=[\frac{\pi}{2}-\delta,\, \frac{\pi}{2}+\delta]$ with
	$0<\delta<\pi/2$ being small and to be chosen. 
	On one hand, thanks to $B\in L^{\infty}$, 
	$K_{21}$ can be controlled simply by 
	\begin{align}\label{es-K21}
		|K_{21}(t,\cdot,\cdot)| \lesssim \delta, 
	\end{align}
	On the other hand, it holds that on the interval $[0,\pi]\backslash I_{\delta},$ 
	\beqs 
	|\cos\theta|\geq \sin\delta\geq \frac{\delta}{2}.
	\eeqs
	Consequently, using the identity $e^{it\sin\theta}=(\frac{\partial_{\theta}}{i t \cos\theta})e^{it\sin\theta}$ and integrating by parts in $\theta,$ we obtain that %for any $\rho\in[0,1],\varphi\in  [0,2\pi]$
	\beq\label{K22}
	\sup_{\varphi\in  [0,2\pi]}  |K_{21}(t,\rho,\varphi)| \lesssim  \frac{1}{\delta t} \sup_{\varphi\in  [0,2\pi]}  
	\big(\|B(\rho \cos (\theta-\varphi))\|_{L_{\theta}^{\infty}}+\|\pa_{\theta}B(\rho \cos (\theta-\varphi))\|_{L_{\theta}^{1}}\big)
	\lesssim  \frac{1}{\delta t}.
	\eeq
	The estimate \eqref{es-K2} follows from 
	%Choose  we find from 
	\eqref{es-K21} and \eqref{K22} by choosing $\delta=Ct^{-\frac12}$. 
\end{proof}

\end{document}
